% =====================================================================
%  Harald Høffding, Benedict Spinoza
%  Translated from the Danish: „Benedikt Spinoza“, Nær og Fjern, vol. 6, no. 233 (17 December 1876), pp. [1]--3.
%
%  TRANSLATION -- GENERATED; DO NOT EDIT BY HAND.  Built by ebook-method/hoeffding-1875-76/tr_parts.py build spinoza-1876 --write
%  from the parts in ebook-method/hoeffding-1875-76/tr/ (GLOSSARY.md, TRANSLATOR-BRIEF.md).  Made from the collated
%  transcription texts/hoeffding/spinoza-1876/transcription.tex (2026-10-15); [N] in the margin = the page of the original.
%  Terms: GLOSSARY.md (choices marked ◆ there for Hans to review).
% =====================================================================
\documentclass[12pt,a4paper,openany]{book}
\usepackage[utf8]{inputenc}
\usepackage[T1]{fontenc}
\usepackage{libertinus}
\usepackage{libertinust1math}
\usepackage[german,french,english]{babel}
\usepackage[protrusion=true,expansion=true]{microtype}
\usepackage{setspace}
\usepackage[margin=1.3in]{geometry}
\usepackage{fancyhdr}
\usepackage{hyperref}
\hypersetup{pdftitle={Benedict Spinoza}, pdfauthor={Harald Høffding}, hidelinks}
\pagestyle{fancy}
\fancyhf{}
\fancyhead[LE,RO]{\thepage}
\fancyhead[C]{\textit{H. Høffding: Benedict Spinoza (1876)}}
\fancypagestyle{plain}{\fancyhf{}\fancyhead[LE,RO]{\thepage}}
\setstretch{1.3}
\usepackage{marginnote}
\usepackage{graphicx}
\DeclareUnicodeCharacter{0254}{\reflectbox{c}}   % ɔ (id est)
\renewcommand{\thefootnote}{\ifcase\value{footnote}\or *\or **\or ***\else\arabic{footnote}\fi)}
\newcommand{\opage}[1]{\marginnote{\footnotesize\textit{[#1]}}}
\newcommand{\signature}[1]{\par\smallskip\hfill #1\hspace*{2em}\par}
\pdfstringdefDisableCommands{\renewcommand{\footnote}[2][]{}\renewcommand{\opage}[1]{}}

\begin{document}

\begin{titlepage}
\begin{center}
{\LARGE\bfseries Near and Far.}\\[0.6em]{\large \emph{A Weekly},}\\[0.2em]{\small edited by P. Hansen.}\\[0.8em]{\small \emph{Sixth Volume} (Nos. 222--274).}\\\vfill{\small Copenhagen 1877. C. A. Reitzel; Otto B. Wroblewsky.}\\[2em]
{\small Translated from the Danish}
\end{center}
\end{titlepage}
\clearpage

% --- p. 1 ---
\setcounter{footnote}{0}%
\opage{1}\begin{center}{\small Copenhagen.  Published every Sunday.\quad {\large\bfseries No. 233.}\quad 17 December 1876.  Subscription: 3 kroner 50 øre a quarter.}\\[0.3em]\rule{\textwidth}{0.4pt}\\[0.8em]{\Large \emph{Benedict Spinoza.}}\\[0.4em]{\small By \textbf{Dr. phil.} Harald Høffding.}\\[0.4em]\rule{3em}{0.4pt}\end{center}\medskip

The zeal of our age for raising statues and memorials
has often brought me to think of the saying about those who
whitewash the tombs of the prophets.  One cannot keep away the melancholy
thought that the recognition of the greatness of the past
does not prevent the misjudgment of contemporaries.  But it is,
on the other hand, natural that a freer and more conciliatory
view comes most readily to the fore in regard to what is past,
where everything is concluded, the struggle ended, and the
position and motives of the combatants illuminated by their place and conditions
within the historical development.  In Spinoza's case
it seems in any event to be so.
Even those who cannot sympathize with his standpoint
and his doctrine cannot deny him their admiration for
his noble and gentle character and must concede his great
significance for later science and literature.  And
when one recalls how much persecution and misjudgment
Spinoza had to suffer, how alone he stood in his time,
and how alien this age felt itself to his ideas, so
that more than 100 years passed before they could break through,
then the statue that is now to be raised to him in The Hague
by contributions not only from his native land but also from
the rest of Europe will stand as a testimony that there has
taken place a progress not only in the understanding of what is great
and significant but also in a liberal-minded and tolerant
way of thinking.

That he stood alone in his time is perhaps somewhat connected
with the fact that he was originally a stranger to the modern
% (column 2)
\setcounter{footnote}{0}%
Western science and literature.  His parents were
Portuguese Jews, who had left their native land for the sake of their
faith and in Holland, the cradle of religious freedom, had
sought a new home.  In Amsterdam there was at that time
a flourishing Jewish congregation with a synagogue and a
rabbinical school, at which renowned teachers worked.  In this city \emph{Benedict Spinoza} was born on 24 November
1632.  On account of the gifts he displayed
early, he was destined for the learned estate, that is to say,
to become a rabbi.  He also made such rapid progress
on this path that already in his sixteenth year he was
a respected Talmudist.  But soon his sharp
understanding and tireless zeal for inquiry began to feel dissatisfied
within the confines of tradition and theology.  Probably
it was at this time that he learned Latin and began
to study the classical authors, the natural sciences
and the Cartesian philosophy.  His teacher in the humanities
was the physician van Ende, who kept a kind of school
in his house and was known for his free religious
views.  A legend tells that Spinoza fell in love
with his daughter, who helped her father with the teaching,
but that she preferred a fellow pupil to him.  Yet in
the most recent time this account has been rejected, since van
Ende's daughter, according to the information that has been found, at
the time of which we are speaking was only a little over ten years
old.

The doubts that arose in him, and the new thoughts
% --- p. 2 ---
\setcounter{footnote}{0}%
\opage{2}which were beginning to take shape, certainly removed him early
somewhat from his co-religionists and made his visits to the synagogue
less frequent; for it is told that they promised
him a not inconsiderable annual pension if he would only visit
the synagogue now and then.  Indeed, it is even said that more
violent means were tried against him.  One evening, as he
was leaving the synagogue, a man attacked him with a dagger, whose
thrust, however, merely slit his clothing.  There
was suspicion that he would go over to Christianity,
and they were loath to yield so gifted an adherent to another
confession.  Yet it must also be emphasized that
the adherents of the Cartesian philosophy at that time were
regarded with suspicion even within the Christian church.
The Second Synod of Dort had recently declared them
excluded from ecclesiastical offices.  The little Jewish congregation,
which had to watch still more strictly over its members,
when neither good nor evil means bore fruit,
finally resorted to the expedient of expelling the apostate from its midst.
This took place in July 1656.  The formula of the ban was
found some years ago in the archive of the Jewish congregation
in Amsterdam; it is composed in the style of the Old Testament
and makes partial use of its words.  ``Cursed,'' --- it says
among other things --- ``be he by day, cursed be he
by night, cursed when he lies down, cursed
when he rises from his bed!  Cursed be his
going out, cursed his coming in!''  ``Know that no one
may have anything to do with him either orally or
in writing, no one may show him any favor, no one dwell
under his roof, no one come near him within less than four
cubits, no one read anything he has written!''

In his twenty-fourth year Spinoza was thus cast out into the world,
% print: stilt (misprint)
violently severed from family and kinsmen, and
without a dwelling anywhere else.  It now stood him in good
stead that the Talmud required everyone who trained to become a
rabbi to learn a trade.  He had learned the
art of grinding optical glasses and had attained great
skill in it.  By this he now procured his frugal
livelihood.  But he was not yet allowed to live in peace.
At the instigation of the Jewish and the Christian clergy
the authorities of Amsterdam banished him in the year 1660 from the city
``to uphold order and subordination.''  He
lived in the following years in Rhynsburg, a small
village near Leiden, where an Arminian sect, the so-called
Collegiants, had found a refuge.  Later he moved
to Voorburg near The Hague, and finally (1669), at the urging of his friends,
to The Hague itself.

Of Spinoza's character and inner life in his youth
one knows only what can be inferred from the outward events
and from the letters and writings of his later years.
His biographers are inclined to infer backwards from the calm and resignation
which are expressed in these, and hold that he also
in his youth displayed such qualities;
yet some hold that one can conceive of them in the
mature thinker only as the result of passionate struggles in
youth.  It is in any event clear that early on he
gathered all his striving and endeavor into one longing ---
that of acquiring a clear and valid knowledge of the nature of man
and his relation to the whole of nature.  His whole
character has an intellectual stamp; investigation and contemplation
are for him the highest thing.  He himself describes in his
treatise on ``the Improvement of the Understanding'' how the
% (column 2)
\setcounter{footnote}{0}%
goods that most people seek, honor, riches and sensual enjoyment,
had lost their value for him, and it now dawned on
him that the true and highest good was ``the knowledge of
the union of our mind with the whole of nature.''  He sought
in the knowledge of truth more than mere theory; he sought
a view of life.  And his greatness consists here precisely in
the harmony that existed in him between knowledge and
character, thought and life.

After the publication of his first writing, an exposition
of the Cartesian philosophy with indications of the points
where it had to be carried further, he soon became famous.
Many sought him out in his dwelling, and he conducted a considerable
correspondence.  One of his correspondents, Oldenburg,
the first President of the Society of Sciences in
London; mentions as the attractive thing about him ``thorough
knowledge combined with gentleness and fine cultivation.''  In the year
1670 he published a ``theological-political treatise,'' which aimed
to show ``that freedom of thought can not only be granted
without injury to piety and to the peace of the state, but that
it cannot even be abolished without these perishing at the same time''.
He demanded here the right for everyone
to interpret the religious doctrines according to his best insight ---
and not only to form his own conviction but also to
express it.  He cites (while seeming to forget
how it had gone for himself) his native city, Amsterdam,
as an example that such freedom can exist with
the nature of an ordered society.  ``In this flourishing and
distinguished city people of all nations and
confessions live in the greatest concord.  There is no sect,
however hated, whose adherents are not (as long as they
injure no one, give everyone his due and
% print: stikkelig (misprint)
live honestly) protected
by the public power of the state''.  He is the first
who has maintained the principle of toleration in this full extent.
The work cited contained moreover historical-critical
investigations of the books of the Old Testament,
whereby the foundation was laid for the whole of modern Biblical criticism,
and rationalistic interpretations of the most important dogmas.
The offense this writing aroused deterred Spinoza
from publishing more.  His chief work, ``Ethics,'' circulated
in copies among his friends and appeared only
after his death.

In the whole last part of his life he suffered from a
consuming chest ailment.  He nevertheless continued his studies steadily,
and with such zeal that on one occasion for three months
he did not leave his house.  He lived with the painter van Spyck,
in whose family he often spent his hours of rest; he liked to talk
especially with his host and his wife about religion
and about the sermons they had heard in church.  Even the day before his death he was down with them to smoke a
pipe of tobacco and talk about the day's sermon.  He constantly
admonished people to remain in their confession and emphasized
the practical-ethical moment in religion.

Spinoza died on 21 February 1677.  No one was present
with him except his friend the physician Ludvig Meyer.
Rumors spread that he was said to have died in
despair or even to have committed suicide, just as
there were also very evil rumors in circulation about his character and conduct.  A Lutheran clergyman, Colerus, who came
to live in the same house in which Spinoza had formerly lived,
undertook to investigate the matter and published a
biography of Spinoza (in Dutch, translated into
% --- p. 3 ---
\setcounter{footnote}{0}%
\opage{3}French: \textit{``La vie de Spinoza, tirée des écrits de ce fameux philosophe et du témoignage de plusieurs personnes dignes de foi, qui l’ont connu particulièrement, par Jean Colerus. A la Haye 1706''} --- the best source-work
for Spinoza's life) and he was able to contradict all
the rumors that prejudice had bred and fanaticism
spread abroad.  Although Colerus himself condemns Spinoza's doctrines,
he emphasizes his pure morals, gentle character
and great unselfishness.  This last quality Spinoza displayed
in particular toward his sisters, to whom he left
his share of the inheritance, just as he several times declined offers of
financial support from his friends.  A professorship at
Heidelberg, which was offered him, he declined so as not to be
hindered in his quiet studies and from fear of intolerance.

One has often, perhaps, derived Spinoza's doctrine exclusively
from Cartesianism.  From a youthful writing found some years ago
it appears that he was strongly influenced
by the Jewish religious philosophers of the Middle Ages and by
Giordano Bruno, the greatest philosopher of the Renaissance, who maintained
the infinity and unity of nature.  What Spinoza
found in Descartes was the clear, mathematical method
and a mechanical conception of nature.  Proceeding from the idea of
the infinite substance, which is at once nature and
God, he seeks to show that it is this which moves in
everything.  The material and the mental world are different
forms for the revelation of the infinite
unity.

Such thoughts had, in the nature of the case, to lie long
uncomprehended; both religious and philosophical prejudices hindered the appropriation of them.  Spinoza lay, according to
Lessing's words, like a dead dog, scorned and kicked by
all parties.  Then came his resurrection at the end of
the eighteenth century with the freer and fresher movement
in poetry, philosophy and science.  It would lead too far
to develop here what he has been for minds such as Lessing,
Herder, Goethe, Novalis, Schelling and Schleiermacher.
For Goethe there breathed in Spinoza's writings an ``air of
peace''; he learned from them to see that there is a great, necessary
connection in everything, and admired the deep resignation
which is expressed in them.  In the writing by Schleiermacher which introduces the religious renewal after the rationalism of the eighteenth
century, he pays homage in an enthusiastic
eulogy to Spinoza's profound and truly religious spirit.
In the lectures that Schelling's disciple Henrik
Steffens held at the beginning of this century in Copenhagen,
and which had so mighty an influence on
the development of our literature, one likewise finds a beautiful
and enthusiastic eulogy of Spinoza.  ``Never can I,'' says
Steffens among other things, ``think of this glorious man in his
calm solitude, with a light in his soul that no accident
could chase away, with an inward peace that
no persecutions could disturb, without a ray of
% (column 2)
\setcounter{footnote}{0}%
his light awakening the most beautiful hope in me.''  (One should
read the whole beginning of the 2nd lecture).

What was attempted at the beginning of the century by the
poetic and speculative path --- to grasp nature and existence
in its wholeness and unity, that is what realistic research is working
toward in our day.  In the natural science
and philosophy of our time there is a strong unity-seeking tendency.  Wide horizons
open, and connections are made and
transitions formed where one earlier halted at a
yawning abyss.  This must naturally revive the memory
of the solitary Dutch thinker of the seventeenth century.
Spinoza is also related to modern
realism in that, in spite of his speculative and deductive method,
he maintains the strict causal connection and demands
a mechanical explanation carried through everywhere.  Especially as regards
the relation between soul and body, his
theory seems to be the one toward which the newest psychological and
physiological investigations point.

We can no longer use Spinoza's method in
philosophy.  What stood for his thought as the first,
the secure starting point: the absolute being, is for us
precisely that which eternally eludes thought.  We must seek
by the way of experience --- which for Spinoza too was the
only one by which we could come to know individual beings
--- to find the laws and forms of the life of the mind and its
connection with the whole of existence.  But Spinozism
stands as a model by its boldness and consistency;
it is an attempt which in this form has its counterpart only in
Hegel's system, and this cannot in originality and in the
purity of its lines compare with Spinoza's.  And besides
the fruitful and profound in the fundamental thought there are
also in Spinoza, as in every great thinker,
subordinate points of much interest and lasting
significance.  I have already mentioned his historical
criticism and his maintaining of the principle of toleration, and shall
only mention further his psychological exposition of the passions,
their origin, development and interplay.  It
is in its kind so classical an exposition that Johannes
Müller as late as some thirty-odd years ago reprinted it
whole and entire in his famous ``Handbuch der Physiologie''.
It has only in our time found counterparts in the expositions of the English
psychologists.

But quite apart from the theoretical results Spinoza stands
as a great model of the character of the true inquirer, and by the gentleness and calm he preserved
in his inner being as well as toward others, although he
had suffered from the conflicts of the view of life, he stands
especially as a figure from whose contemplation one can, as Steffens said, draw ``the most beautiful hope'' in a time such as ours, in which views clash no less, even if
the clash perhaps takes place under less violent
forms.

\end{document}
